\section{Three Different Approaches to Deep Causal Estimation}
\label{section:mainbody}

The architectures proposed in the deep learning literature for causal estimation build upon the core idea discussed above. First, we introduce ``S-Learners" and ``T-Learners" to show how neural networks can be used to estimate non-linearities in potential outcomes. Second, given the right objectives, a neural network can learn representations of the treated and control distributions that are deconfounded (Fig. \ref{fig:balancing}). This approach, which can be related theoretically to non-parametric matching, is illustrated by the foundational TARNet algorithm in section \ref{section:ipw} \citep{Shalit2017}. Finally, the machine learning for causal inference literature has been largely driven by the introduction of semi-parametric frameworks that allow predictive machine learning models to be plugged-in to doubly robust estimation equations \citep{van2011targeted,chernozhukov2016double,chernozhukov2021automatic}. In section \ref{section:ipw}, we introduce the concept of influence functions and the targeted maximum likelihood estimator to explain the Dragonnet algorithm. For clarity the algorithms presented here all share a familial resemblence to the TARNet algorithm. However, we note that there are many other approaches to using deep learning for causal inference (e.g., the generative models described in Appendix A.\ref{appendix:generative}).
\iffalse
\textcolor{blue}{(There is a new approach using invariant risk minimization (IRM) for causal estimation. Paper: Invariant Representation Learning for Treatment Effect Estimation)}
\begin{table}[]
\caption{Summary of Potential Outcome Approaches. Columns are selection on observable estimation strategies, rows are machine learning estimation techniques. Note that papers may appear multiple times.}
\label{tab:review}
\resizebox{\textwidth}{!}{%
\begin{tabular}{@{}llll@{}}
\toprule
 & Balancing & IPW & Other \\ \midrule
Integral Probability Metric & \citet{Johansson2016,Shalit2017}\\
                            & \citet{Johansson2018,Du2019} & \citet{Johansson2018}  \\ \cmidrule(r){1-1}
Propensity-Matching & \citet{Yao2018,Schwab2018} &  \\ \cmidrule(r){1-1}
Adversarial Training & \citet{Johansson2018,Du2019} & \citet{KallusDeepMatch,Ozery-flato}\\
& & \citet{Johansson2018} & \citet{Yoon2018} \\ \cmidrule(r){1-1}
Semiparametric Metaloss &  & \citet{Shi2019} &  \\ \cmidrule(r){1-1}
Selection Bias Dropout & \citet{Alaa2017} &  \\ \cmidrule(r){1-1}
Reconstruction Loss & \citet{Yao2018,Du2019} & \citet{Atan2018} \\ \cmidrule(r){1-1}
\bottomrule
\end{tabular}%
}
\end{table}
\fi